\setcounter{chapter}{7}
\chapter{Lebesgue積分}

\paragraph{本章の目標}
可測関数が準備できたので，いよいよLebesgue積分を定義する．
単関数，非負可測関数，符号付可測関数の順に測度空間上の積分を定義し，
その基本性質とRiemann積分との関係を確認する．

\section{単関数の積分}

以後，本章の一般論では測度空間 $(X,\calM,\mu)$ を固定する．

\begin{definition}[非負単関数]\label{def:nonnegative-simple-function}
写像 $s:X \to [0,\infty)$ が非負の単関数（simple function）または階段関数（step function）であるとは，
\textbf{有限個}の非負実数 $a_1,\dots,a_n$ と
互いに素な可測集合 $A_1,\dots,A_n \subset X$ を用いて
$s=\sum_{k=1}^n a_k 1_{A_k}$ と書けることをいう．
\end{definition}

\begin{remark}
単関数は関数値が有限個しかない可測関数である．
%積分の最初のモデルとして最も計算しやすい．
特に指示関数 $1_A$ は最も基本的な単関数である．
\end{remark}

\begin{definition}[単関数の積分]\label{def:simple-function-integral}
非負単関数 $s=\sum_{k=1}^n a_k 1_{A_k}$ に対して，
\[
\int_X s\,d\mu
:=
\sum_{k=1}^n a_k\,\mu(A_k)
\]
と定める．
\end{definition}

\begin{proposition}[単関数の積分のwell-definedness]\label{prop:simple-function-integral-well-defined}
非負単関数 $s$ の表示の仕方によらず，$\int_X s\,d\mu$ は同じ値を与える．
\end{proposition}
\begin{proof}
$s=\sum_{i=1}^n a_i 1_{A_i}=\sum_{j=1}^m b_j 1_{B_j}$ と2通りに表したとする．
ここで $\{A_i\}$ と $\{B_j\}$ はそれぞれ互いに素で可測である．
各 $i,j$ に対して $C_{ij}:=A_i \cap B_j$ とおくと，$\{C_{ij}\}_{i,j}$ は互いに素な可測集合族であり，
$A_i=\bigsqcup_{j=1}^m C_{ij}$，$B_j=\bigsqcup_{i=1}^n C_{ij}$ が成り立つ．
また $x \in C_{ij}$ なら $s(x)=a_i=b_j$ だから，$C_{ij}$ 上では係数が一致する．
ゆえに測度の可算加法性より
\[
\sum_{i=1}^n a_i\,\mu(A_i)
=\sum_{i,j}a_i\,\mu(C_{ij})
=\sum_{i,j}b_j\,\mu(C_{ij})
=\sum_{j=1}^m b_j\,\mu(B_j).
\]
したがって積分値は表示に依らない．
\end{proof}

\begin{example}
\begin{enumerate}
    \item 可測集合 $A \subset X$ に対して $\int_X 1_A\,d\mu = \mu(A)$ である
    \item 定数関数 $c \ge 0$ に対して $\int_X c\,d\mu = c\,\mu(X)$ である
    \item $s=\sum_{k=1}^n a_k 1_{A_k}$ なら，積分は各値 $a_k$ にその値をとる領域の測度 $\mu(A_k)$ を掛けて足したものである
\end{enumerate}
\end{example}

\begin{proposition}[単関数の基本性質]\label{prop:simple-function-integral-basic}
非負単関数 $s,t$ と $c \ge 0$ に対して次が成り立つ．
\begin{enumerate}
    \item $\int_X (s+t)\,d\mu = \int_X s\,d\mu + \int_X t\,d\mu$
    \item $\int_X cs\,d\mu = c \int_X s\,d\mu$
    \item $s \le t$ なら $\int_X s\,d\mu \le \int_X t\,d\mu$
\end{enumerate}
\end{proposition}
\begin{proof}
まず $s=\sum_{i=1}^n a_i 1_{A_i}$，$t=\sum_{j=1}^m b_j 1_{B_j}$ と書く．
共通細分 $C_{ij}:=A_i \cap B_j$ を取れば，$s+t=\sum_{i=1}^n \sum_{j=1}^m (a_i+b_j)1_{C_{ij}}$ である．
したがって
\[
\int_X (s+t)\,d\mu
=\sum_{i,j}(a_i+b_j)\mu(C_{ij})
=\sum_i a_i \mu(A_i)+\sum_j b_j \mu(B_j)
=\int_X s\,d\mu+\int_X t\,d\mu.
\]
これで (1) が従う．

(2) は $cs=\sum_{i=1}^n (ca_i)1_{A_i}$ と書けば直ちに従う．

(3) を示す．
$t-s$ は非負単関数であるから (1), (2) より
$\int_X t\,d\mu=\int_X s\,d\mu+\int_X (t-s)\,d\mu\ge \int_X s\,d\mu$ である．
\end{proof}

\section{非負可測関数の積分}

\begin{definition}[非負可測関数の積分]\label{def:nonnegative-lebesgue-integral}
非負可測関数 $f:X \to [0,\infty]$ に対して
\[
\int_X f\,d\mu
:=
\sup\left\{\int_X s\,d\mu \;\middle|\; 0 \le s \le f,\ s \text{ は非負単関数}\right\}
\]
と定める．
\end{definition}

\begin{remark}
単関数は値が有限個しかないので計算しやすい．
したがって上の定義は，
「$f$ の下側から入る単関数でどこまで積分値を近づけられるか」
を表している．
つまり，複雑な可測関数を
より簡単な単関数で近似して積分している．
\end{remark}

\begin{lemma}[単関数による近似]\label{lem:simple-function-approximation}
非負可測関数 $f:X \to [0,\infty]$ に対して，
単関数列 $\{s_n\}_{n=1}^{\infty}$ で $0 \le s_1 \le s_2 \le \cdots \le f$ かつ
$s_n(x) \to f(x)$ $(x \in X)$ を満たすものが存在する．
\end{lemma}
\begin{proof}
各 $n \in \NN$ に対して
\[
s_n(x)
:=
\begin{cases}
\dfrac{k}{2^n}
&\left(
\dfrac{k}{2^n} \le f(x) < \dfrac{k+1}{2^n},
\ \ 0 \le k \le n2^n-1
\right),\\[1.2ex]
n & (f(x)\ge n)
\end{cases}
\]
と定める．
すると $s_n$ は有限個の値しかとらず，
各値をとる集合は $\{x \in X \mid k/2^n \le f(x) < (k+1)/2^n\}$ または
$\{x \in X \mid f(x)\ge n\}$ の形で書けるので可測である．
したがって $s_n$ は単関数である．

また定義から $0 \le s_n(x) \le f(x)$ が成り立つ．
さらに，$f(x)<\infty$ のとき $n>f(x)$ なら $0 \le f(x)-s_n(x)<2^{-n}$ であり，
$f(x)=\infty$ のときは $s_n(x)=n \to \infty$ である．
したがって $s_n(x)\to f(x)$ である．

最後に，二進展開の精度が上がるので $s_n(x)\le s_{n+1}(x)$ が成り立つ．
よって所望の単関数列を得る．
\end{proof}

\begin{theorem}[非負可測関数の積分は単関数近似の極限]\label{thm:nonnegative-integral-simple-approximation}
$f:X \to [0,\infty]$ を非負可測関数とし，
$s_n$ を $0 \le s_1 \le s_2 \le \cdots \le f$ かつ $s_n \to f$ を満たす単関数列とする．
このとき $\int_X f\,d\mu=\lim_{n\to\infty}\int_X s_n\,d\mu$ が成り立つ．
\end{theorem}
\begin{proof}
各 $n$ について $s_n \le f$ だから，
定義より $\int_X s_n\,d\mu \le \int_X f\,d\mu$ である．
また $s_n \le s_{n+1}$ だから，単関数の単調性より
$\int_X s_n\,d\mu \le \int_X s_{n+1}\,d\mu$ である．
したがって列 $\{\int_X s_n\,d\mu\}_{n=1}^{\infty}$ は単調増加であり，
$L:=\lim_{n\to\infty}\int_X s_n\,d\mu$ が存在して $L \le \int_X f\,d\mu$ である．

逆向きの不等式を示すため，
$t \le f$ を満たす任意の非負単関数 $t=\sum_{j=1}^m a_j 1_{A_j}$ を取る．
ここで $a_j \ge 0$，$A_j$ は互いに素な可測集合とする．
各 $j,n$ に対して
$A_{j,n}:=A_j \cap \{x \in X \mid s_n(x)\ge a_j-1/n\}$ とおく．
$s_n(x)\to f(x)\ge a_j$ が $A_j$ 上で成り立つので，
$A_{j,1}\subset A_{j,2}\subset \cdots$ かつ $\bigcup_{n=1}^{\infty}A_{j,n}=A_j$ である．

そこで
$u_n:=\sum_{j=1}^m (a_j-1/n)_+ 1_{A_{j,n}}$ とおくと，$u_n$ は非負単関数であり，
$A_{j,n}$ 上では $(a_j-1/n)_+ \le s_n$ だから $u_n \le s_n$ が成り立つ．
したがって
$\int_X u_n\,d\mu \le \int_X s_n\,d\mu \le L$ である．

一方，測度の下からの連続性より
$\mu(A_{j,n}) \to \mu(A_j)$ であるから，
\[
\int_X u_n\,d\mu
=\sum_{j=1}^m \left(a_j-\frac1n\right)_+ \mu(A_{j,n})
\to \sum_{j=1}^m a_j \mu(A_j)
=\int_X t\,d\mu.
\]
よって
$\int_X t\,d\mu \le L$ である．
$t \le f$ を満たす単関数 $t$ は任意であったから，
定義より $\int_X f\,d\mu \le L$ である．以上より $\int_X f\,d\mu=L$ が従う．
\end{proof}

\begin{proposition}[非負可測関数の基本性質]\label{prop:nonnegative-integral-basic}
非負可測関数 $f,g:X \to [0,\infty]$ と $c \ge 0$ に対して次が成り立つ．
\begin{enumerate}
    \item $\int_X (f+g)\,d\mu = \int_X f\,d\mu + \int_X g\,d\mu$
    \item $\int_X cf\,d\mu = c\int_X f\,d\mu$
    \item $f \le g$ なら $\int_X f\,d\mu \le \int_X g\,d\mu$
\end{enumerate}
\end{proposition}
\begin{proof}
単関数列 $s_n \uparrow f$，$t_n \uparrow g$ を取る．
すると $s_n+t_n \uparrow f+g$，$cs_n \uparrow cf$ である．
したがって\cref{thm:nonnegative-integral-simple-approximation,prop:simple-function-integral-basic}より
\[
\int_X (f+g)\,d\mu
=\lim_{n\to\infty}\int_X (s_n+t_n)\,d\mu
=\lim_{n\to\infty}\left(\int_X s_n\,d\mu+\int_X t_n\,d\mu\right)
=\int_X f\,d\mu+\int_X g\,d\mu,
\]
および
$\int_X cf\,d\mu=\lim_{n\to\infty}\int_X cs_n\,d\mu
=\lim_{n\to\infty} c\int_X s_n\,d\mu=c\int_X f\,d\mu$
を得る．

また $f\le g$ なら，$f$ 以下の非負単関数はすべて $g$ 以下でもある．
したがって上限をとれば $\int_X f\,d\mu\le \int_X g\,d\mu$ であり，(3) が従う．
\end{proof}

\begin{definition}[部分集合上の積分]\label{def:integral-over-subset}
$A \subset X$ を可測集合とする．
非負可測関数 $f:X \to [0,\infty]$ に対して
$\int_A f\,d\mu := \int_X f1_A\,d\mu$ と定める．
符号付可測関数の場合も同様に定める．
\end{definition}

\begin{proposition}[集合に関する加法性]\label{prop:nonnegative-integral-disjoint-additivity}
$A,B \subset X$ を互いに素な可測集合とし，
$f:X \to [0,\infty]$ を非負可測関数とする．
このとき $\int_{A \sqcup B} f\,d\mu=\int_A f\,d\mu+\int_B f\,d\mu$ が成り立つ．
\end{proposition}
\begin{proof}
$1_{A \sqcup B}=1_A+1_B$ だから，非負可測関数の加法性より
\[
\int_{A \sqcup B} f\,d\mu
=\int_X f1_{A \sqcup B}\,d\mu
=\int_X f1_A\,d\mu+\int_X f1_B\,d\mu
=\int_A f\,d\mu+\int_B f\,d\mu.
\]
\end{proof}

\section{零集合と almost everywhere}

\begin{proposition}[零集合上の積分]\label{prop:null-set-integral-nonnegative}
$N \subset X$ が可測で $\mu(N)=0$ ならば，
任意の非負可測関数 $f:X \to [0,\infty]$ に対して
$\int_N f\,d\mu=0$ が成り立つ．
\end{proposition}
\begin{proof}
定義より
\[
\int_N f\,d\mu
=\sup\left\{\int_N s\,d\mu \;\middle|\; 0 \le s \le f,\ s \text{ は単関数}\right\}.
\]
ところが $N$ 上の任意の単関数 $s=\sum_{k=1}^n a_k1_{A_k}$ について，
各 $A_k \subset N$ だから $\mu(A_k)\le \mu(N)=0$ である．
したがって $\int_N s\,d\mu=\sum_{k=1}^n a_k \mu(A_k)=0$ である．
ゆえに上限も $0$ である．
\end{proof}

\begin{proposition}[非負関数の a.e. 不変性]\label{prop:nonnegative-integral-ae-invariance}
$f,g:X \to [0,\infty]$ を非負可測関数とする．
$f=g$ a.e.\ on $X$ ならば
$\int_X f\,d\mu=\int_X g\,d\mu$ が成り立つ．
\end{proposition}
\begin{proof}
$N:=\{x \in X \mid f(x)\ne g(x)\}$ とおくと $\mu(N)=0$ である．
したがって\cref{prop:null-set-integral-nonnegative}より
$\int_N f\,d\mu=\int_N g\,d\mu=0$ である．
また $X=(X\setminus N)\sqcup N$ だから，
集合に関する加法性より
\[
\int_X f\,d\mu
=\int_{X\setminus N} f\,d\mu+\int_N f\,d\mu
=\int_{X\setminus N} g\,d\mu+\int_N g\,d\mu
=\int_X g\,d\mu.
\]
\end{proof}

\begin{proposition}[非負関数の積分が $0$ であることの判定]\label{prop:nonnegative-integral-zero-iff-ae-zero}
非負可測関数 $f:X \to [0,\infty]$ に対して，次は同値である．
\begin{enumerate}
    \item $\int_X f\,d\mu=0$
    \item $f=0$ a.e.\ on $X$
\end{enumerate}
\end{proposition}
\begin{proof}
\textbf{(2) $\Rightarrow$ (1)}\quad
$f=0$ a.e.\ なら，\cref{prop:nonnegative-integral-ae-invariance}を $g=0$ に適用して
$\int_X f\,d\mu=\int_X 0\,d\mu=0$ である．

\textbf{(1) $\Rightarrow$ (2)}\quad
各 $n \in \NN$ に対して
$A_n:=\{x \in X \mid f(x)\ge 1/n\}$ とおく．
すると
$n^{-1}1_{A_n}\le f$ だから単調性より
\[
\frac1n \mu(A_n)
=\int_X \frac1n 1_{A_n}\,d\mu
\le \int_X f\,d\mu=0.
\]
よって $\mu(A_n)=0$ である．
一方，$\{x \in X \mid f(x)>0\}=\bigcup_{n=1}^{\infty}A_n$ だから可算加法性より $\mu(\{f>0\})=0$ である．
したがって $f=0$ a.e.\ である．
\end{proof}

\section{符号付可測関数の積分}

\begin{definition}[正部分と負部分]\label{def:positive-negative-parts}
可測関数 $f:X \to \eRR$ に対して
$f^+:=\max(f,0)$，$f^-:=\max(-f,0)$ を $f$ の正部分，負部分という．
\end{definition}

\begin{remark}
常に $f=f^+-f^-$，$|f|=f^++f^-$ が成り立つ．
したがって符号付関数の積分は，
非負関数の積分に帰着できる．
\end{remark}

\begin{definition}[可積分関数]\label{def:lebesgue-integrable-function}
可測関数 $f:X \to \eRR$ が可積分であるとは，
$\int_X f^+\,d\mu<\infty$ かつ $\int_X f^-\,d\mu<\infty$ が成り立つことをいう．
このとき
$\int_X f\,d\mu:=\int_X f^+\,d\mu-\int_X f^-\,d\mu$ で定める．
\end{definition}

\begin{remark}
$\infty-\infty$ は定義しないので，
正部分と負部分の両方が有限であるときにだけ
符号付関数の積分を定義する．
この条件は，後で見るように $\int_X |f|\,d\mu<\infty$ と同値である．
\end{remark}

\begin{proposition}[表示の独立性]\label{prop:integral-representation-independence}
$f=u-v=u'-v'$ と書け，
$u,v,u',v'$ がすべて非負可測で
$\int_X u\,d\mu,\ \int_X v\,d\mu,\ \int_X u'\,d\mu,\ \int_X v'\,d\mu<\infty$
とする．このとき
$\int_X u\,d\mu-\int_X v\,d\mu=\int_X u'\,d\mu-\int_X v'\,d\mu$ が成り立つ．
\end{proposition}
\begin{proof}
$u-v=u'-v'$ だから $u+v'=u'+v$ である．
両辺は非負可測関数なので，加法性より
\[
\int_X u\,d\mu+\int_X v'\,d\mu
=\int_X (u+v')\,d\mu
=\int_X (u'+v)\,d\mu
=\int_X u'\,d\mu+\int_X v\,d\mu.
\]
したがって移項して
$\int_X u\,d\mu-\int_X v\,d\mu
=\int_X u'\,d\mu-\int_X v'\,d\mu$ を得る．
\end{proof}

\begin{theorem}[可積分性と絶対値]\label{thm:integrable-iff-absolute}
可測関数 $f:X \to \eRR$ に対して，次は同値である．
\begin{enumerate}
    \item $f$ は可積分である
    \item $|f|$ は可積分である
\end{enumerate}
さらに，このとき
$\int_X |f|\,d\mu=\int_X f^+\,d\mu+\int_X f^-\,d\mu$ が成り立つ．
\end{theorem}
\begin{proof}
$|f|=f^++f^-$ だから，非負可測関数の加法性より
$\int_X |f|\,d\mu=\int_X f^+\,d\mu+\int_X f^-\,d\mu$ である．
したがって右辺が有限であることと
$\int_X f^+\,d\mu<\infty$，$\int_X f^-\,d\mu<\infty$ が成り立つことは同値である．
これはまさに $f$ の可積分性の定義である．
\end{proof}

\begin{corollary}[優関数による可積分性の判定]\label{cor:dominated-integrability}
$f:X \to \eRR$ を可測関数とし，
$g:X \to [0,\infty]$ を可積分な非負可測関数とする．
もし $|f|\le g$ a.e.\ on $X$ ならば $f$ は可積分である．
\end{corollary}
\begin{proof}
$N:=\{x \in X \mid |f(x)|>g(x)\}$ とおくと $\mu(N)=0$ である．
したがって
$|f|=|f|1_{X\setminus N}$ a.e.\ であり，非負関数の a.e.\ 不変性より
$\int_X |f|\,d\mu=\int_X |f|1_{X\setminus N}\,d\mu$ である．
しかも $X\setminus N$ 上では $|f|1_{X\setminus N}\le g$ だから単調性より
$\int_X |f|\,d\mu\le \int_X g\,d\mu<\infty$ である．
ゆえに\cref{thm:integrable-iff-absolute}より $f$ は可積分である．
\end{proof}

\begin{corollary}[可積分関数は有限値を a.e.\ でとる]\label{cor:integrable-finite-ae}
$f:X \to \eRR$ が可積分ならば
$|f|<\infty$ a.e.\ on $X$ が成り立つ．
\end{corollary}
\begin{proof}
$N:=\{x \in X \mid |f(x)|=\infty\}$ とおく．
各 $n \in \NN$ について
$n1_N\le |f|$ だから，単調性より
$n\,\mu(N)=\int_X n1_N\,d\mu \le \int_X |f|\,d\mu<\infty$ である．
これがすべての $n$ で成り立つためには
$\mu(N)=0$ でなければならない．
\end{proof}

\begin{proposition}[零集合上の任意の関数]\label{prop:null-set-integral-signed}
$N \subset X$ が可測で $\mu(N)=0$ ならば，
任意の可測関数 $f:X \to \eRR$ は $N$ 上で可積分であり，
$\int_N f\,d\mu=0$ が成り立つ．
\end{proposition}
\begin{proof}
非負可測関数については零集合上の積分は $0$ である
（\cref{prop:null-set-integral-nonnegative}）．
したがって $\int_N f^+\,d\mu=\int_N f^-\,d\mu=0$ である．
よって $f$ は $N$ 上で可積分で，
$\int_N f\,d\mu=\int_N f^+\,d\mu-\int_N f^-\,d\mu=0$ となる．
\end{proof}

\begin{proposition}[符号付関数の a.e.\ 不変性]\label{prop:signed-integral-ae-invariance}
$f,g:X \to \eRR$ を可測関数とする．
$f=g$ a.e.\ on $X$ で，$f$ が可積分ならば $g$ も可積分であり，
$\int_X f\,d\mu=\int_X g\,d\mu$ が成り立つ．
\end{proposition}
\begin{proof}
$N:=\{x \in X \mid f(x)\ne g(x)\}$ とおくと $\mu(N)=0$ である．
$X\setminus N$ 上では $|g|=|f|$ であり，
$N$ 上の積分はどちらも $0$ だから
\[
\int_X |g|\,d\mu
=\int_{X\setminus N}|g|\,d\mu+\int_N |g|\,d\mu
=\int_{X\setminus N}|f|\,d\mu
\le \int_X |f|\,d\mu<\infty.
\]
よって $g$ も可積分である．

さらに
\[
\int_X f\,d\mu
=\int_{X\setminus N} f\,d\mu+\int_N f\,d\mu
=\int_{X\setminus N} g\,d\mu+\int_N g\,d\mu
=\int_X g\,d\mu.
\]
\end{proof}

\section{積分の絶対連続性}

\begin{theorem}[積分の絶対連続性]\label{thm:integral-absolute-continuity}
$f:X \to \eRR$ を可積分関数とする．
このとき任意の $\eps>0$ に対して，ある $\delta>0$ が存在して，
可測集合 $A \subset X$ が $\mu(A)<\delta$ を満たせば
$\int_A |f|\,d\mu<\eps$ が成り立つ．
\end{theorem}
\begin{proof}
$h:=|f|$ とおくと，$h$ は非負可積分である．
積分の定義より，ある非負単関数 $s \le h$ が存在して
$\int_X h\,d\mu-\int_X s\,d\mu<\eps/2$ となる．
いま $s=\sum_{k=1}^n a_k1_{A_k}$ と書き，$M:=\max_{1\le k\le n}a_k$ とおく．

$M=0$ なら $s=0$ だから
$\int_X h\,d\mu<\eps/2$ である．
この場合は任意の $\delta>0$ でよい．

$M>0$ の場合は $\delta:=\eps/(2M)$ とおく．
可測集合 $A \subset X$ が $\mu(A)<\delta$ を満たすとする．
すると集合に関する加法性と $s\le h$ から
\[
\int_X h\,d\mu
=\int_A h\,d\mu+\int_{X\setminus A}h\,d\mu
\ge \int_A h\,d\mu+\int_{X\setminus A}s\,d\mu.
\]
よって
\[
\int_A h\,d\mu
\le \int_X h\,d\mu-\int_{X\setminus A}s\,d\mu
=\left(\int_X h\,d\mu-\int_X s\,d\mu\right)+\int_A s\,d\mu
<\frac{\eps}{2}+M\,\mu(A)
<\frac{\eps}{2}+M\delta=\eps.
\]
これで示された．
\end{proof}

\section{具体例}

\begin{example}[Dirichlet関数]
$f:=1_{\QQ \cap [0,1]}$ とする．このとき $\int_{[0,1]} f\,d\lambda=0$ である．
\end{example}
\begin{proof}
\cref{cor:countable-lebesgue-null}より $\lambda(\QQ \cap [0,1])=0$ である．
したがって $\int_{[0,1]} f\,d\lambda=\int_{[0,1]} 1_{\QQ \cap [0,1]}\,d\lambda=\lambda(\QQ \cap [0,1])=0$ である．
\end{proof}

\begin{remark}
Dirichlet関数はRiemann積分できなかったが，Lebesgue積分では $0$ と自然に積分できる．
ここに，Lebesgue積分が
「どこで値をとるか」と「その領域の面積」を直接結びつけていることが表れている．
\end{remark}

\begin{example}
$g:=1_{[0,1]\setminus \QQ}$ とすると $g=1$ a.e.\ on $[0,1]$ だから
$\int_{[0,1]} g\,d\lambda=1$ である．
\end{example}
\begin{proof}
$N:=\QQ \cap [0,1]$ は零集合であり，$[0,1]\setminus N$ 上で $g=1_{[0,1]\setminus N}=1$ である．
したがって a.e.\ 不変性より $\int_{[0,1]} g\,d\lambda=\int_{[0,1]} 1\,d\lambda=\lambda([0,1])=1$ である．
\end{proof}

\section{Riemann積分との関係}

\begin{theorem}[Riemann可積分ならLebesgue可積分]\label{thm:riemann-implies-lebesgue-integrable}
$I \subset \RR^d$ を有界閉区間とし，
$f:I \to \RR$ をRiemann可積分関数とする．
このとき $f$ はLebesgue可測かつLebesgue可積分であり，$\int_I f\,d\lambda$ はRiemann積分値に一致する．
\end{theorem}
\begin{proof}
Riemann可積分性より，任意の $n \in \NN$ に対して
分割 $\calP_n$ が存在して $U(f,\calP_n)-L(f,\calP_n)<1/n$ となる．
この分割の各小区間を $R_{n,1},\dots,R_{n,N_n}$ と書き，
各小区間上の下限と上限を
$\ell_{n,j}:=\inf_{x \in R_{n,j}} f(x)$，$u_{n,j}:=\sup_{x \in R_{n,j}} f(x)$ とおく．
すると $\phi_n:=\sum_{j=1}^{N_n} \ell_{n,j}1_{R_{n,j}}$，
$\psi_n:=\sum_{j=1}^{N_n} u_{n,j}1_{R_{n,j}}$ は単関数で，$\phi_n \le f \le \psi_n$ が成り立つ．
さらに
$\int_I \phi_n\,d\lambda = L(f,\calP_n)$，$\int_I \psi_n\,d\lambda = U(f,\calP_n)$ である．

ここで $g:=\sup_{n\in\NN}\phi_n$，$h:=\inf_{n\in\NN}\psi_n$ とおく．
単関数は可測だから $g,h$ は可測であり，$g \le f \le h$ である．
しかも各 $n$ に対して $0 \le h-g \le \psi_n-\phi_n$ だから単調性より
\[
0 \le \int_I (h-g)\,d\lambda
\le \int_I (\psi_n-\phi_n)\,d\lambda
=U(f,\calP_n)-L(f,\calP_n)
<\frac1n.
\]
したがって $\int_I (h-g)\,d\lambda=0$ であり，非負関数の積分が $0$ であることの判定より
$h-g=0$ a.e.\ on $I$ である．
特に $N:=\{x \in I \mid g(x)<h(x)\}$ は零集合であり，$I\setminus N$ 上では $f=g=h$ である．

よって $f$ は可測であることを示す．
実数 $a$ に対して
\[
\{x \in I \mid f(x)>a\}
=\bigl(\{g>a\}\cap(I\setminus N)\bigr)\cup\bigl(\{f>a\}\cap N\bigr)
\]
と書ける．
$g$ は可測，$N$ は零集合だから可測であり，
しかもLebesgue測度の完備性（\cref{thm:lebesgue-null-measurable}）より
$N$ の任意の部分集合は可測である．
したがって右辺は可測である．
ゆえに $f$ はLebesgue可測である．

次に可積分性を示す．
Riemann可積分関数は有界だから，ある $M>0$ が存在して
$|f(x)|\le M$ $(x \in I)$ である．したがって $|f|\le M1_I$ である．
$I$ は有界閉区間なので
$\int_I M1_I\,d\lambda = M\,\lambda(I)=M|I|<\infty$ である．
よって\cref{cor:dominated-integrability}より $f$ はLebesgue可積分である．

最後に積分値の一致を示す．
各 $n$ に対して $\phi_n \le f \le \psi_n$ だから単調性より
\[
L(f,\calP_n)=\int_I \phi_n\,d\lambda
\le \int_I f\,d\lambda
\le \int_I \psi_n\,d\lambda=U(f,\calP_n).
\]
Riemann可積分性より $U(f,\calP_n)-L(f,\calP_n)\to 0$ である．
したがってはさみうちにより
$\int_I f\,d\lambda=\lim_{n\to\infty}L(f,\calP_n)=\lim_{n\to\infty}U(f,\calP_n)$ である．
これはRiemann積分値そのものである．
\end{proof}

\begin{remark}
この定理により，Lebesgue積分はRiemann積分の拡張になっている．
ただし本質は，単に対象が増えただけではなく，
零集合と極限に関してはるかに安定な理論になっている点にある．
\end{remark}

\begin{example}[広義Riemann積分可能だがLebesgue積分可能でない関数]
$f(x):=\sin x/x$ $(x \ge 1)$ を考える．
このとき $\int_1^{\infty} \frac{\sin x}{x}\,dx$ は広義Riemann積分として収束するが，
$f$ はLebesgue可積分ではない．
\end{example}
\begin{proof}
まず広義Riemann積分の収束を示す．
$1<R<\infty$ に対して部分積分を行うと
\[
\int_1^R \frac{\sin x}{x}\,dx
=\left[-\frac{\cos x}{x}\right]_{x=1}^{x=R}
-\int_1^R \frac{\cos x}{x^2}\,dx.
\]
右辺第1項は $R\to\infty$ で極限をもち，
第2項は
\[
\int_1^{\infty}\frac{|\cos x|}{x^2}\,dx
\le \int_1^{\infty}\frac{1}{x^2}\,dx
<\infty
\]
だから収束する．
よって $\int_1^{\infty} \frac{\sin x}{x}\,dx$ は広義Riemann積分として収束する．

次にLebesgue可積分でないことを示す．
各 $k \in \NN$ に対して
$I_k:=[k\pi+\pi/6,\ k\pi+5\pi/6]$ とおくと，$x \in I_k$ では $|\sin x|\ge 1/2$ である．
したがって
\[
\int_1^{\infty}\frac{|\sin x|}{x}\,dx
\ge \sum_{k=1}^{\infty}\int_{I_k}\frac{|\sin x|}{x}\,dx
\ge \frac12\sum_{k=1}^{\infty}\int_{I_k}\frac{1}{x}\,dx
\ge \frac{\pi}{3}\sum_{k=1}^{\infty}\frac{1}{k\pi+\frac{5\pi}{6}}.
\]
右辺は調和級数と同程度に発散するから
$\int_1^{\infty}\frac{|\sin x|}{x}\,dx=\infty$ である．
よって $\sin x/x$ はLebesgue可積分ではない．
\end{proof}

\section{解釈とまとめ}

測度空間上の積分では，まず $\int 1_A\,d\mu = \mu(A)$ という形で「集合の測度」を積分に埋め込み，
そこから単関数の積分を定め，
さらに一般の可測関数へ拡張した．
したがってこの積分の本質は，
関数値そのものよりも「その値をとる領域の大きさ」を正確に足し合わせることにある．

特に本章で得た重要な事実は次の通りである．
\begin{enumerate}
    \item Dirichlet関数はLebesgue積分でき，積分値は $0$ である
    \item $f$ が可積分であることと $|f|$ が可積分であることは同値である
    \item 可積分関数は零集合上で値を変えても積分値が変わらない
    \item Riemann可積分関数はすべてLebesgue可積分である
\end{enumerate}

次章では，この積分を実際に計算で運用するための基本公式
$\int(f+g)$，$\int cf$，$\int |f|$，$\int_X f$ を体系的に整理する．
